\begin{theorem}[Rank of diagonal coefficients across a periodic cut]
    \label{thm:mpo_diagonal_cut_rank}
    \lean{MPOTensor.diagonalCutMatrix,
      MPOTensor.diagonalCutMatrix_rank_le}
    \leanok
    \uses{def:mpo_tensor}
    Split a periodic chain into consecutive blocks of lengths $L$ and $R$.
    For words $x\in\{0,\ldots,d-1\}^L$ and
    $y\in\{0,\ldots,d-1\}^R$, set
    \begin{align}
      P_{x,y}
      &=
      \tr\!\left(M^{x_0x_0}\cdots M^{x_{L-1}x_{L-1}}
        M^{y_0y_0}\cdots M^{y_{R-1}y_{R-1}}\right).
      \notag
    \end{align}
    Then the ordinary complex rank of $P$ satisfies
    $\rank_{\C}P\le D^2$.
\end{theorem}
\begin{proof}\leanok
    \uses{lem:mpo_eval_word_append}
    For each word, write
    $A_x=M^{x_0x_0}\cdots M^{x_{L-1}x_{L-1}}$ and
    $B_y=M^{y_0y_0}\cdots M^{y_{R-1}y_{R-1}}$. The identity
    \begin{align}
      P_{x,y}
      &=\tr(A_xB_y)
       =\sum_{a,b=0}^{D-1}(A_x)_{ab}(B_y)_{ba}
      \notag
    \end{align}
    factors $P$ through the $D^2$-dimensional space indexed by pairs $(a,b)$.
\end{proof}

\begin{theorem}[Rank under nonzero scalar multiplication]
    \label{thm:matrix_rank_smul}
    \lean{Matrix.rank_smul_of_ne_zero}
    \leanok
    Let $K$ be a field, let $A$ be a finite matrix over $K$, and let
    $a\in K$ be nonzero. Then $\rank_{K}(aA)=\rank_{K}A$.
\end{theorem}
\begin{proof}\leanok
    Multiplication by $a$ identifies the column spaces, with inverse given by
    multiplication by $a^{-1}$.
\end{proof}

\begin{theorem}[Rank under faithful scalar extension]
    \label{thm:matrix_rank_scalar_extension}
    \lean{Matrix.rank_map_algebraMap}
    \leanok
    Let $K\subseteq L$ be a faithful field extension, let
    $\iota:K\hookrightarrow L$ be the inclusion, and let $A$ be a finite
    matrix over $K$. Then
    $\rank_{L}\!\left((\iota(A_{ij}))_{ij}\right)=\rank_{K}A$.
\end{theorem}
\begin{proof}\leanok
    For every finite selection of columns, the selected columns are linearly
    independent over $K$ if and only if their scalar extensions are linearly
    independent over $L$. Selecting a basis from each column space proves the
    two rank inequalities.
\end{proof}

\begin{definition}[Joint probability distribution and marginals]
    \label{def:finite_joint_distribution_marginals}
    \lean{Matrix.IsJointDistribution,
      Matrix.rowMarginal,
      Matrix.columnMarginal}
    \leanok
    Let $X$ and $Y$ be finite sets. A matrix
    $P=(P_{x,y})_{x\in X,y\in Y}$ is a joint probability distribution if
    $P_{x,y}\geq0$ for every $x\in X$ and $y\in Y$, and
    \begin{align}
      \sum_{x\in X}\sum_{y\in Y}P_{x,y}
      &=1.
      \notag
    \end{align}
    Its row and column marginals are respectively
    \begin{align}
      p_x&=\sum_{y\in Y}P_{x,y},
      \notag\\
      q_y&=\sum_{x\in X}P_{x,y}.
      \notag
    \end{align}
\end{definition}

\begin{definition}[Entropy and classical mutual information]
    \label{def:finite_classical_mutual_information}
    \lean{Entropy.probabilityEntropy,
      Entropy.classicalMutualInformation}
    \leanok
    \uses{def:finite_joint_distribution_marginals}
    For $t>0$, set $h(t)=-t\log t$, and set $h(0)=0$. The entropy of a
    probability distribution $a=(a_z)_{z\in Z}$ on a finite set is
    $H(a)=\sum_{z\in Z}h(a_z)$. For a joint probability distribution $P$
    with row and column marginals $p$ and $q$, put
    \begin{align}
      H(P)&=\sum_{x\in X}\sum_{y\in Y}h(P_{x,y}),
      \notag\\
      I(X:Y)_P&=H(p)+H(q)-H(P).
      \notag
    \end{align}
\end{definition}

\begin{lemma}[Entropy is bounded by the logarithm of the support cardinality]
    \label{lem:probability_entropy_le_log_card_support}
    \lean{Entropy.probabilityEntropy_le_log_card_support}
    \leanok
    \uses{def:finite_classical_mutual_information}
    Let $p=(p_z)_{z\in Z}$ be a probability distribution on a finite set $Z$.
    Define $h(t)=-t\log t$ for $t>0$ and $h(0)=0$. Then
    \begin{align}
      \sum_{z\in Z}h(p_z)
      &\leq
      \log\left\lvert\{z\in Z:p_z\neq0\}\right\rvert.
      \notag
    \end{align}
\end{lemma}
\begin{proof}\leanok
    Let $k=\lvert\{z\in Z:p_z\neq0\}\rvert$. Jensen's inequality for the
    concave function $h$, with uniform weights on the support, gives
    \begin{align}
      \frac{1}{k}\sum_{\substack{z\in Z\\p_z\neq0}}h(p_z)
      &\leq
      h\left(\frac{1}{k}
        \sum_{\substack{z\in Z\\p_z\neq0}}p_z\right)
       =h\left(\frac{1}{k}\right)
       =\frac{\log k}{k}.
      \notag
    \end{align}
    Multiplication by $k$ proves the result, since values outside the support
    contribute $h(0)=0$.
\end{proof}

\begin{theorem}[Classical mutual information is bounded by ordinary rank]
    \label{thm:classical_mutual_information_le_log_rank}
    \lean{Entropy.classicalMutualInformation_le_log_rank}
    \leanok
    \uses{def:finite_classical_mutual_information}
    Let $X$ and $Y$ be finite sets and let $P=(P_{x,y})_{x\in X,y\in Y}$
    be a joint probability distribution. With natural logarithms,
    $I(X:Y)_P\leq\log\rank_{\R}P$. Equivalently, with logarithms to base two,
    $2^{I(X:Y)_P}\leq\rank_{\R}P$.
\end{theorem}
\begin{proof}\leanok
    \uses{lem:probability_entropy_le_log_card_support}
    This is Theorem~4.1 of \cite{Riazanov2017PSDRank}. If the number of
    nonzero rows exceeds the ordinary real rank, choose a nontrivial real
    linear relation among those rows. Nonnegativity forces the relation to
    have coefficients of both signs. Rescaling each row by
    $1+\varepsilon\beta_x$ gives two endpoint distributions. For every
    $y\in Y$, the row relation gives
    \begin{align}
      \sum_x(1+\varepsilon\beta_x)P_{x,y}
      &=\sum_xP_{x,y}
        +\varepsilon\underbrace{\sum_x\beta_xP_{x,y}}_{=0}
       =\sum_xP_{x,y}.
      \notag
    \end{align}
    Thus each endpoint preserves every column marginal, removes at least one
    row, and has no larger rank.
    Write $h(t)=-t\log t$ for $t>0$, with $h(0)=0$, and put
    $m=\min_x\beta_x<0<M=\max_x\beta_x$. Then
    \begin{align}
      I(X:Y)_{P^{(\varepsilon)}}
      &=I(X:Y)_P+\varepsilon S(\beta),
      \notag\\
      S(\beta)
      &=\sum_x\beta_x\left(h(p_x)-\sum_y h(P_{x,y})\right).
      \notag
    \end{align}
    If $S(\beta)\geq0$, choose $\varepsilon=-m^{-1}\geq0$; otherwise choose
    $\varepsilon=-M^{-1}\leq0$. In either case
    $\varepsilon S(\beta)\geq0$, so the selected endpoint has mutual
    information no smaller than the original distribution.
    Repeating the construction leaves at most
    $\rank_{\R}P$ nonzero rows. Finally,
    $I(X:Y)\leq H(X)$, and the entropy of a distribution supported on $k$
    points is at most $\log k$.
\end{proof}

\begin{theorem}[Classical estimate for normalized diagonal cut coefficients]
    \label{thm:mpo_diagonal_cut_classical_mutual_information}
    \lean{MPOTensor.diagonalCut_classicalMutualInformation_le_two_log}
    \leanok
    \uses{thm:mpo_diagonal_cut_rank,
      thm:classical_mutual_information_le_log_rank,
      thm:matrix_rank_smul,
      thm:matrix_rank_scalar_extension}
    Let $D\geq1$, and let $W$ be a real diagonal coefficient matrix across a
    periodic cut of an MPO with bond dimension $D$. Suppose its scalar
    extension to $\C$ is the cut matrix $P^{\C}$ from
    Theorem~\ref{thm:mpo_diagonal_cut_rank}.
    If $Z>0$ and $P=Z^{-1}W$ is a joint probability distribution, then
    $I(X:Y)_P\leq2\log D$.
\end{theorem}
\begin{proof}\leanok
    Normalization does not change real rank by
    Theorem~\ref{thm:matrix_rank_smul}. Faithful scalar extension does
    not change rank by Theorem~\ref{thm:matrix_rank_scalar_extension}; hence
    Theorem~\ref{thm:mpo_diagonal_cut_rank} gives
    $\rank_{\R}P\leq D^2$.
    Theorem~\ref{thm:classical_mutual_information_le_log_rank} gives
    $I(X:Y)_P\leq\log\rank_{\R}P$, and monotonicity of the
    logarithm gives the result.
\end{proof}

\begin{definition}[Normalized diagonal distribution at fixed length]
    \label{def:mpo_normalized_diagonal_cut_distribution}
    \lean{MPOTensor.IsDiagonalAt,
      MPOTensor.IsDiagonal,
      MPOTensor.diagonalCutRealMatrix,
      MPOTensor.diagonalCutMass,
      MPOTensor.normalizedDiagonalCutDistribution}
    \leanok
    \uses{def:mpo_tensor}
    For a chain of length $N=L+R$, call $\rho^{(N)}(M)$ diagonal if
    $\rho^{(N)}(M)_{u,v}=0$ whenever $u\ne v$. Call $M$ globally diagonal if
    $\rho^{(N)}(M)$ is diagonal for every $N>0$. For configurations $x$ and
    $y$ on the two consecutive blocks, define
    \begin{align}
      W_{x,y}&=\re\rho^{(L+R)}(M)_{(x,y),(x,y)},
      \notag\\
      Z&=\sum_{x,y}W_{x,y},
      \notag\\
      \widehat P&=Z^{-1}W.
      \notag
    \end{align}
\end{definition}

\begin{lemma}[Diagonal coefficient mass and trace]
    \label{lem:mpo_diagonal_cut_mass_trace}
    \lean{MPOTensor.diagonalCutMass_eq_trace_re}
    \leanok
    \uses{def:mpo_normalized_diagonal_cut_distribution}
    For every MPO tensor $M$ and cut lengths $L,R$, the total diagonal
    coefficient mass satisfies $Z=\re\tr\rho^{(L+R)}(M)$.
\end{lemma}
\begin{proof}\leanok
    Pairing each configuration $\sigma$ on $L+R$ sites with its unique block
    split $(x,y)$ on $L$ and $R$ sites gives
    \begin{align}
      Z
      &=\sum_{x,y}W_{x,y}
       =\sum_\sigma\re\rho^{(L+R)}(M)_{\sigma,\sigma}
       =\re\tr\rho^{(L+R)}(M).
      \notag
    \end{align}
\end{proof}

\begin{theorem}[Classical estimate for a positive diagonal finite-chain operator]
    \label{thm:mpo_diagonal_finite_chain_classical_mutual_information}
    \lean{MPOTensor.diagonalCutMatrix_apply_eq_mpo,
      MPOTensor.diagonalCutRealMatrix_map_ofReal_of_posSemidef,
      MPOTensor.diagonalCutRealMatrix_nonneg_of_posSemidef,
      MPOTensor.normalizedDiagonalCutDistribution_isJointDistribution,
      MPOTensor.diagonalFiniteChain_classicalMutualInformation_le_two_log}
    \leanok
    \uses{def:mpo_normalized_diagonal_cut_distribution,
      thm:mpo_diagonal_cut_rank}
    Let $D\geq1$. Suppose that $\rho^{(L+R)}(M)$ is diagonal and positive
    semidefinite, and that $Z>0$. Then $W_{x,y}\geq0$, the scalar extension of
    $W$ to $\C$ is the diagonal cut matrix of
    Theorem~\ref{thm:mpo_diagonal_cut_rank}, and $\widehat P$ is a joint
    probability distribution. Its classical mutual information satisfies
    $I(X:Y)_{\widehat P}\leq2\log D$.
\end{theorem}
\begin{proof}\leanok
    \uses{thm:mpo_diagonal_cut_classical_mutual_information}
    Positive semidefiniteness makes every diagonal entry real and non-negative,
    so $W_{x,y}\geq0$ and its scalar extension recovers the complex diagonal
    cut matrix. Since $Z$ is the sum of the entries of $W$,
    $\sum_{x,y}\widehat P_{x,y}
      =Z^{-1}\sum_{x,y}W_{x,y}=1$.
    Thus $\widehat P$ is a joint probability distribution. Diagonality says
    that these coefficients constitute the full classical finite-chain state,
    and Theorem~\ref{thm:mpo_diagonal_cut_classical_mutual_information} gives
    the bound.
\end{proof}

\begin{corollary}[Classical estimate for a diagonal matrix product density operator]
    \label{cor:mpo_diagonal_classical_mutual_information}
    \lean{MPOTensor.diagonalMPDO_classicalMutualInformation_le_two_log}
    \leanok
    Let $D\geq1$, and let $M$ generate diagonal positive semidefinite operators
    at every positive chain length. If $L+R>0$ and the trace at length $L+R$
    is positive, then the classical mutual information of its normalized
    diagonal coefficients satisfies $I(X:Y)\leq2\log D$.
\end{corollary}
\begin{proof}\leanok
    \uses{thm:mpo_diagonal_finite_chain_classical_mutual_information,
      lem:mpo_diagonal_cut_mass_trace}
    Global diagonality and positivity give the fixed-length hypotheses of
    Theorem~\ref{thm:mpo_diagonal_finite_chain_classical_mutual_information}.
    Since the chain operator is diagonal,
    \begin{align}
      Z
      &=\sum_{x,y}W_{x,y}
       =\re\tr\rho^{(L+R)}(M),
      \notag
    \end{align}
    so the trace positivity hypothesis supplies the required normalization.
\end{proof}

\begin{definition}[Operator-Schmidt rank]
    \label{def:bipartite_operator_schmidt_rank}
    \lean{Matrix.HasOperatorSchmidtDecomposition,
      Matrix.operatorSchmidtRank, Matrix.operatorSchmidtRank_minimal}
    \leanok
    For a bipartite density operator
    $\rho\in\MN{d_A}\otimes\MN{d_B}$, its
    \emph{operator-Schmidt rank} is the least integer $r$ for which there are
    matrices $A_t\in\MN{d_A}$ and $B_t\in\MN{d_B}$ satisfying
    \begin{align}
        \rho&=\sum_{t=1}^{r}A_t\otimes B_t.
        \notag
    \end{align}
    No Hermiticity or positivity condition is imposed on the factors. This is
    the definition in \cite[Equation~(1)]{DeLasCuevas2019Separability}; the
    same formula defines the rank of an arbitrary complex bipartite matrix.
\end{definition}

\begin{theorem}[Two virtual boundaries bound the operator-Schmidt rank]
    \label{thm:mpdo_periodic_cut_operator_schmidt_rank}
    \lean{MPOTensor.periodicCutLeftFactor,
      MPOTensor.periodicCutRightFactor,
      MPOTensor.bipartitionedNormalizedMPO_hasOperatorSchmidtDecomposition,
      MPOTensor.operatorSchmidtRank_bipartitionedNormalizedMPO_le}
    \leanok
    \uses{def:mpdo_split_normalized_mpo,
      def:bipartite_operator_schmidt_rank}
    Let $M$ be an MPO tensor of operator bond dimension $D$, and let
    $N=L+K$. Opening the two virtual bonds at the consecutive cut gives
    \begin{align}
        \sigma^{(N)}_{L:K}
        &=\sum_{a,b=0}^{D-1} X_{ab}\otimes Y_{ab},
        \notag
    \end{align}
    and
    \begin{align}
        \operatorname{OSR}(\sigma^{(N)}_{L:K})&\leq D^2.
        \notag
    \end{align}
    For left-block words $x,x'$ and right-block words $y,y'$, the factors are
    \begin{align}
        (X_{ab})_{x,x'}
        &=\tr[\rho^{(N)}(M)]^{-1}
          \bigl(M^{x_0x'_0}\cdots M^{x_{L-1}x'_{L-1}}\bigr)_{ab},
        \notag
    \end{align}
    and
    \begin{align}
        (Y_{ab})_{y,y'}
        &=\bigl(M^{y_0y'_0}\cdots M^{y_{K-1}y'_{K-1}}\bigr)_{ba}.
        \notag
    \end{align}
    This is the identity
    $\tr(AB)=\sum_{a,b}A_{ab}B_{ba}$, with the normalization scalar
    absorbed into the left factor. The algebraic decomposition does not use
    positivity. Since $0^{-1}=0$, it also remains valid when the trace
    vanishes; a nonzero trace is needed only to interpret
    $\sigma^{(N)}_{L:K}$ as the normalized physical state.

    This is the two-virtual-boundary algebraic step toward the finite
    mutual-information estimate in
    \cite[Proposition~4.5]{Cirac2016MPDO_arXiv}. It does not by itself prove
    $I_L\leq4\log D$: such a conclusion would require an entropy bound in
    terms of the operator-Schmidt rank.
\end{theorem}

\begin{proof}\leanok
    Split each closed virtual word into its left and right products. Expanding
    the trace over the two exposed virtual indices gives the displayed sum of
    $D^2$ product matrices. The minimality property of the operator-Schmidt
    rank then bounds it by the number of displayed terms.
\end{proof}


\begin{theorem}[Coefficient-space characterization of operator-Schmidt rank]
    \label{thm:operator_schmidt_coefficient_space}
    \lean{Matrix.range_operatorSchmidtMap_eq_operatorBlockSpan,
      Matrix.operatorSchmidtRank_eq_finrank_operatorBlockSpan,
      Matrix.operatorSchmidtRank_minimal}
    \leanok
    \uses{def:bipartite_operator_schmidt_rank}
    Write $\rho_{ij}\in\MN{d_B}$ for the blocks determined by a basis
    of the first factor, and define
    $\mathcal R_\rho(X)=\sum_{i,j}X_{ij}\rho_{ij}$. Then
    \begin{align}
        \operatorname{OSR}(\rho)
        &=\dim\range\mathcal R_\rho
         =\dim\spn\{\rho_{ij}:1\leq i,j\leq d_A\}.
        \notag
    \end{align}
\end{theorem}
\begin{proof}\leanok
    The range of $\mathcal R_\rho$ is the span of the blocks. If
    $\rho=\sum_{t=1}^{r}A_t\otimes B_t$, every block belongs to
    $\spn\{B_1,\ldots,B_r\}$, so the range dimension is at most
    $r$. Conversely, choose a basis $B_1,\ldots,B_s$ of the block span and
    expand each block as $\rho_{ij}=\sum_t(A_t)_{ij}B_t$. This gives
    $\rho=\sum_{t=1}^{s}A_t\otimes B_t$, where
    $s=\dim\range\mathcal R_\rho$.
\end{proof}

\begin{theorem}[Operator-Schmidt rank under local isometries]
    \label{thm:operator_schmidt_rank_local_isometry}
    \lean{Matrix.operatorSchmidtRank_local_mul_le,
      Matrix.operatorSchmidtRank_local_isometry_conj}
    \leanok
    \uses{def:bipartite_operator_schmidt_rank}
    Let $V_A:\mathcal H_A\to\mathcal K_A$ and
    $V_B:\mathcal H_B\to\mathcal K_B$ be isometries between finite-dimensional
    complex spaces. For every operator $X$ on
    $\mathcal H_A\otimes\mathcal H_B$,
    \begin{align}
      \operatorname{OSR}\!\left(
        (V_A\otimes V_B)X(V_A\otimes V_B)^\dagger\right)
      &=\operatorname{OSR}(X).
      \notag
    \end{align}
    This remains valid when one or more of the spaces have dimension zero.
\end{theorem}
\begin{proof}\leanok
    \uses{thm:operator_schmidt_coefficient_space}
    More generally, multiplying on the left and right by product matrices
    cannot increase operator-Schmidt rank: apply the four local matrices to
    the two factors in each term of a shortest product decomposition. Applying
    this inequality first to $V_A,V_B$ and then to their adjoints gives the two
    inequalities. The identities $V_A^\dagger V_A=1$ and
    $V_B^\dagger V_B=1$ identify the twice-transformed operator with $X$.
\end{proof}

\begin{theorem}[Compression to both marginal supports]
    \label{thm:operator_schmidt_rank_marginal_support_compression}
    \lean{Matrix.PosSemidef.eq_marginalSupport_expansion_compression,
      Matrix.PosSemidef.operatorSchmidtRank_marginalSupport_compression}
    \leanok
    \uses{def:bipartite_operator_schmidt_rank}
    Let $\rho\geq0$ be an operator on $H_A\otimes H_B$. Let
    $V_A:\widehat H_A\to H_A$ and $V_B:\widehat H_B\to H_B$ be isometries
    whose range projectors are the support projectors $P_A$ and $P_B$ of the
    two marginals. Set $W=V_A\otimes V_B$ and
    $\rho_c=W^\dagger\rho W$. Then
    \begin{align}
      W\rho_cW^\dagger&=\rho,
      &\operatorname{OSR}(\rho_c)&=\operatorname{OSR}(\rho).
      \notag
    \end{align}
    No marginal is required to be faithful, and zero-dimensional coordinate
    spaces are permitted.
\end{theorem}
\begin{proof}\leanok
    \uses{thm:marginal_support_left_absorption,
      thm:marginal_support_right_absorption,
      thm:right_marginal_support_left_absorption,
      thm:right_marginal_support_right_absorption,
      thm:operator_schmidt_rank_local_isometry}
    The four marginal-support absorption identities give
    \begin{align}
      (P_A\otimes P_B)\rho&=\rho,
      &\rho(P_A\otimes P_B)&=\rho.
      \notag
    \end{align}
    Since $WW^\dagger=P_A\otimes P_B$, associativity yields
    \begin{align}
      W\rho_cW^\dagger
      &=(WW^\dagger)\rho(WW^\dagger)=\rho.
      \notag
    \end{align}
    The local-isometry theorem applied to $\rho_c$ gives
    $\operatorname{OSR}(W\rho_cW^\dagger)=\operatorname{OSR}(\rho_c)$.
    Substitution of the reconstruction identity proves the rank equality.
\end{proof}

\begin{lemma}[Partial traces under product isometries]
    \label{lem:partial_trace_product_isometry}
    \lean{Matrix.partialTraceRight_kronecker_conj_of_right_isometry,
      Matrix.partialTraceLeft_kronecker_conj_of_left_isometry}
    \leanok
    \uses{def:partial_trace, def:partial_trace_left}
    Let $A:H_A\to K_A$ and $B:H_B\to K_B$ be linear maps between
    finite-dimensional complex spaces, and let $X$ be an operator on
    $H_A\otimes H_B$. If $B^\dagger B=1$, then
    \begin{align}
      \tr_{K_B}\!\left((A\otimes B)X(A\otimes B)^\dagger\right)
      &=A(\tr_{H_B}X)A^\dagger.
      \notag
    \end{align}
    If $A^\dagger A=1$, then, symmetrically,
    \begin{align}
      \tr_{K_A}\!\left((A\otimes B)X(A\otimes B)^\dagger\right)
      &=B(\tr_{H_A}X)B^\dagger.
      \notag
    \end{align}
    The identities include zero-dimensional spaces.
\end{lemma}
\begin{proof}\leanok
    Expand the matrix entries in orthonormal bases. In the first identity, the
    sum over the traced output basis gives the matrix entries of
    $B^\dagger B=1$, and hence collapses the two input indices on $H_B$.
    The second identity follows by the same argument with the factors
    interchanged.
\end{proof}

\begin{theorem}[Faithful marginals after simultaneous support compression]
    \label{thm:faithful_marginals_support_compression}
    \lean{Matrix.PosSemidef.partialTraceRight_marginalSupport_compression,
      Matrix.PosSemidef.partialTraceLeft_marginalSupport_compression,
      Matrix.PosSemidef.marginalSupport_compression_marginals_posDef}
    \leanok
    \uses{thm:operator_schmidt_rank_marginal_support_compression}
    In the setting of
    Theorem~\ref{thm:operator_schmidt_rank_marginal_support_compression}, the
    two marginals of $\rho_c$ satisfy
    \begin{align}
      \tr_{\widehat H_B}\rho_c
      &=V_A^\dagger(\tr_{H_B}\rho)V_A,
      &
      \tr_{\widehat H_A}\rho_c
      &=V_B^\dagger(\tr_{H_A}\rho)V_B.
      \notag
    \end{align}
    Both compressed marginals are positive definite, including when one of
    their spaces has dimension zero.
\end{theorem}
\begin{proof}\leanok
    \uses{lem:partial_trace_product_isometry,
      thm:compression_on_support_posDef,
      thm:operator_schmidt_rank_marginal_support_compression}
    Insert the reconstruction $\rho=W\rho_cW^\dagger$ into the two covariance
    identities of Lemma~\ref{lem:partial_trace_product_isometry}, and multiply
    by the corresponding adjoint isometries. Each displayed marginal is the
    compression of a positive semidefinite matrix to its support, so its
    positive definiteness follows from
    Theorem~\ref{thm:compression_on_support_posDef}.
\end{proof}

\begin{theorem}[Channel of a faithful bipartite marginal]
    \label{thm:supported_marginal_channel}
    \lean{Matrix.finrank_range_supportedMarginalMap,
      Matrix.supportedMarginalMap_isKrausCPTP,
      Matrix.supportedMarginalMap_diagonal,
      Matrix.supportedMarginalReconstruction_eq}
    \leanok
    \uses{def:bipartite_operator_schmidt_rank,
      thm:operator_schmidt_coefficient_space}
    Let $\rho\geq0$ be a bipartite complex matrix whose first marginal is
    faithful. Choose an eigenbasis in which
    $\tr_B\rho=\diag(p_1,\ldots,p_{d_A})$ with
    every $p_i>0$, and define the linear map $\Phi_\rho$ on matrix units by
    \begin{align}
        \Phi_\rho(E_{ij})
        &=\frac{\rho_{ij}}{\sqrt{p_ip_j}}.
        \notag
    \end{align}
    Write $\sigma=\diag(p_1,\ldots,p_{d_A})$ and
    $\tau=\tr_A\rho$. Then $\Phi_\rho$ is completely positive and trace
    preserving,
    \begin{align}
        \Phi_\rho(\sigma)&=\tau,
        &\dim\range\Phi_\rho&=\operatorname{OSR}(\rho),
        \notag
    \end{align}
    and application of
    $\Phi_\rho$ to the first half of
    $\sum_{i,j}\sqrt{p_ip_j}\,E_{ij}\otimes E_{ij}$, followed by restoring the
    order of the two factors, reconstructs $\rho$.
\end{theorem}
\begin{proof}\leanok
    Strict positivity of the $p_i$ makes the entrywise input scaling
    invertible, so it does not change the range. Positivity of $\rho$ gives a
    Kraus representation of $\Phi_\rho$, while the first marginal equation
    gives trace preservation. Summing the diagonal blocks gives
    $\Phi_\rho(\sigma)=\tau$. Substitution on matrix units proves the
    reconstruction identity. This is the finite-dimensional Choi
    representation \cite{Choi1975CompletelyPositive} with the input marginal
    absorbed into the canonical purification.
\end{proof}

\section{Support compression for entropy functionals}

% Issues #5229, #5241, and #5245 establish one- and two-marginal support
% compression, including preservation of operator-Schmidt rank and
% faithfulness of both compressed marginals.  The remaining analytic and
% integration steps are tracked in #5211.

\begin{definition}[Isometric conjugation into a larger matrix algebra]
    \label{def:isometric_nonunital_conjugation}
    \lean{Matrix.isometryConjNonUnitalStarAlgHom,
      Matrix.isometryConjNonUnitalStarAlgHom_apply}
    \leanok
    Let $V:\C^k\to\C^D$ be an isometry, so that
    $V^\dagger V=\Id_k$. Define $\iota_V:\MN{k}\to\MN{D}$ by
    \begin{align}
      \iota_V(A)&=VAV^\dagger.
      \label{eq:mpdo_isometric_conjugation}
    \end{align}
    The map $\iota_V$ is complex-linear, multiplicative, and $*$-preserving.
    In general it is not unital:
    $\iota_V(\Id_k)=VV^\dagger$ is the projection onto the range of $V$.
\end{definition}

\begin{theorem}[Functional calculus under rectangular isometric conjugation]
    \label{thm:cfc_rectangular_isometry_zero}
    \lean{Matrix.cfc_conj_isometry_of_zero}
    \leanok
    Let $A\in\MN{k}$ be Hermitian, let $V:\C^k\to\C^D$ satisfy
    $V^\dagger V=\Id_k$, and let $f:\R\to\R$ satisfy $f(0)=0$. Then
    \begin{align}
      f(VAV^\dagger)&=Vf(A)V^\dagger,
      \label{eq:mpdo_cfc_isometric_conjugation}
    \end{align}
    where both sides use the continuous functional calculus on the respective
    finite spectra.
\end{theorem}

\begin{proof}\leanok
    \uses{def:isometric_nonunital_conjugation}
    Apply functoriality of the continuous functional calculus to
    $\iota_V$. The condition $f(0)=0$ removes the contribution from the
    orthogonal complement of the range of $V$, where $\iota_V(A)$ vanishes.
\end{proof}

\begin{theorem}[Trace preservation under isometric inclusion]
    \label{thm:isometric_embedding_trace}
    \lean{Matrix.trace_mul_mul_conjTranspose_of_conjTranspose_mul_eq_one}
    \leanok
    Let $V:\C^k\to\C^D$ satisfy $V^\dagger V=\Id_k$. For every
    $A\in\MN{k}$,
    \begin{align}
      \tr(VAV^\dagger)&=\tr(A).
      \label{eq:mpdo_isometric_trace}
    \end{align}
\end{theorem}

\begin{proof}\leanok
    Cyclicity of the trace gives
    $\tr(VAV^\dagger)=\tr(V^\dagger VA)=\tr(A)$.
\end{proof}

\begin{theorem}[Compression cancels isometric expansion]
    \label{thm:isometric_compression_expansion}
    \lean{Matrix.conjTranspose_mul_mul_mul_conjTranspose_mul_of_isometry}
    \leanok
    Let $V:\C^k\to\C^D$ satisfy $V^\dagger V=\Id_k$. For every
    $A\in\MN{k}$,
    \begin{align}
      V^\dagger(VAV^\dagger)V&=A.
      \label{eq:mpdo_isometric_compression_expansion}
    \end{align}
\end{theorem}

\begin{proof}\leanok
    Associativity and $V^\dagger V=\Id_k$ reduce the left-hand side to
    $(V^\dagger V)A(V^\dagger V)=A$.
\end{proof}

\begin{theorem}[Reconstruction from a reference support]
    \label{thm:support_compression_reconstruction}
    \lean{Matrix.PosSemidef.eq_expansion_compression_of_kernel_le_of_mul_conjTranspose_eq_supportProj}
    \leanok
    Let $\rho,\omega\in\MN{D}$ be positive semidefinite and suppose that
    $\ker\omega\subseteq\ker\rho$. If
    $V:\C^k\to\C^D$ satisfies $VV^\dagger=P_{\supp(\omega)}$, then
    \begin{align}
      V(V^\dagger\rho V)V^\dagger&=\rho.
      \label{eq:mpdo_support_reconstruction}
    \end{align}
    No condition on $V^\dagger V$ is needed for this identity.
\end{theorem}

\begin{proof}\leanok
    The kernel inclusion and positivity give
    $\rho P_{\supp(\omega)}=\rho$; taking adjoints also gives
    $P_{\supp(\omega)}\rho=\rho$. Substituting
    $VV^\dagger=P_{\supp(\omega)}$ on both sides of $\rho$ proves
    (\ref{eq:mpdo_support_reconstruction}).
\end{proof}

\begin{theorem}[Relative entropy under support compression]
    \label{thm:relative_entropy_support_compression}
    \lean{TNLean.quantumRelativeEntropy_support_compression}
    \leanok
    \uses{def:quantum_relative_entropy}
    Let $\rho,\omega\in\MN{D}$ be positive semidefinite with
    $\ker\omega\subseteq\ker\rho$. Let $V:\C^k\to\C^D$ satisfy
    \begin{align}
      V^\dagger V&=\Id_k,\\
      VV^\dagger&=P_{\supp(\omega)}.
      \label{eq:mpdo_reference_support_isometry}
    \end{align}
    Then
    \begin{align}
      D(\rho\Vert\omega)
      &=D(V^\dagger\rho V\Vert V^\dagger\omega V).
      \label{eq:mpdo_relative_entropy_support_compression}
    \end{align}
    The logarithm is totalized by $\log 0=0$ on the kernel.
\end{theorem}

\begin{proof}\leanok
    \uses{def:isometric_nonunital_conjugation,
      thm:cfc_rectangular_isometry_zero,
      thm:isometric_embedding_trace,
      thm:support_compression_reconstruction}
    Theorem~\ref{thm:support_compression_reconstruction} reconstructs both
    $\rho$ and $\omega$ from their compressions. Apply
    Theorem~\ref{thm:cfc_rectangular_isometry_zero} to the logarithm, use
    multiplicativity of $\iota_V$, and finish with
    (\ref{eq:mpdo_isometric_trace}).
\end{proof}

\begin{theorem}[Functional calculus under transpose]
    \label{thm:cfc_transpose_hermitian}
    \lean{Matrix.cfc_transpose}
    \lean{Matrix.IsHermitian.log_transpose}
    \lean{Matrix.PosSemidef.rpow_transpose}
    \lean{Matrix.PosSemidef.sqrt_transpose}
    \lean{Matrix.IsHermitian.supportProj_transpose}
    \leanok
    \uses{def:hermitian_support_projection}
    If $A\in\MN{n}$ is Hermitian and $f:\R\to\R$, then
    \begin{align}
      f(A)^T&=f(A^T).
      \label{eq:mpdo_cfc_transpose}
    \end{align}
    In particular,
    \begin{align}
      \log(A^T)&=(\log A)^T, &
      P_{A^T}&=P_A^T.
      \label{eq:mpdo_log_support_transpose}
    \end{align}
    If $A$ is positive semidefinite, then for every $r\in\R$,
    \begin{align}
      (A^T)^r&=(A^r)^T, &
      \sqrt{A^T}&=(\sqrt A)^T.
      \label{eq:mpdo_power_sqrt_transpose}
    \end{align}
    These identities include singular matrices and matrices whose index set is
    empty. They are project-derived rather than statements of CPSV16.
\end{theorem}

\begin{proof}\leanok
    \uses{def:hermitian_support_projection}
    For a Hermitian matrix, transpose is entrywise complex conjugation.
    Apply covariance of the continuous functional calculus under this real
    star-algebra automorphism. Continuity of $f$ is needed only on the finite
    spectrum of $A$, where it is automatic. The logarithm, real-power, and
    square-root identities follow by specialization. For the support
    projection, specialize to the function that is one away from zero and zero
    at zero.
\end{proof}

\begin{theorem}[Logarithm of a tensor product on its support]
    \label{thm:mpdo_support_log_kronecker}
    \lean{Matrix.log_kronecker_posSemidef}
    \leanok
    \uses{def:hermitian_support_projection}
    Let $A$ and $B$ be positive semidefinite, with $P_A$ and $P_B$ the
    orthogonal projections onto their respective ranges. Then
    \begin{align}
      \log(A\otimes B)
      &=(\log A)\otimes P_B+P_A\otimes(\log B).
      \label{eq:mpdo_support_log_tensor_product}
    \end{align}
    Here the logarithm is extended by zero on the kernel. Neither matrix is
    required to be positive definite, and either index set may be empty.
    This auxiliary identity is project-derived rather than a theorem of
    CPSV16.
\end{theorem}

\begin{proof}\leanok
    \uses{def:hermitian_support_projection, lem:cfc_conj_unitary}
    Diagonalize $A$ and $B$. On a product eigenvector with eigenvalues
    $a,b\geq0$, equation~(\ref{eq:mpdo_support_log_tensor_product}) becomes
    \begin{align}
      \log(ab)
      &=\log(a)\mathbf 1_{b\ne0}
        +\mathbf 1_{a\ne0}\log(b).
      \notag
    \end{align}
    If $a,b>0$, this is the ordinary product identity for the logarithm. If
    either eigenvalue is zero, both sides vanish. Conjugating the resulting
    diagonal identity by the product eigenbasis proves the formula.
\end{proof}

\begin{theorem}[Logarithm of a positive square root]
    \label{thm:log_sqrt_possemidef}
    \lean{Matrix.PosSemidef.cfc_log_sqrt}
    \leanok
    If $\rho\succeq0$, then
    \begin{align}
      \log\sqrt{\rho}&=\frac12\log\rho.
      \label{eq:mpdo_log_sqrt}
    \end{align}
\end{theorem}

\begin{proof}\leanok
    On every positive eigenvalue this is the scalar identity
    $\log\sqrt{x}=\frac12\log x$. At $x=0$, both sides vanish under the
    convention $\log0=0$.
\end{proof}

\begin{theorem}[Logarithm of a real power]
    \label{thm:log_rpow_posdef}
    \lean{Matrix.PosDef.cfc_log_rpow}
    \leanok
    If $A\succ0$, then, for every $s\in\R$,
    \begin{align}
      \log(A^s)&=s\log A.
      \label{eq:mpdo_log_rpow}
    \end{align}
\end{theorem}

\begin{proof}\leanok
    Every eigenvalue of $A$ is positive, so the assertion follows from
    $\log(x^s)=s\log x$ on $(0,\infty)$ and the continuous functional
    calculus.
\end{proof}

\begin{theorem}[Support-aware vector-state Jensen inequality for the logarithm]
    \label{thm:support_log_vector_jensen}
    \lean{Matrix.PosSemidef.re_dotProduct_cfc_log_mulVec_le_log}
    \leanok
    \uses{def:hermitian_support_projection}
    Let $A\in\MN{n}$ be positive semidefinite and let $v\in\C^n$ satisfy
    $\langle v,v\rangle=1$ and $P_{\supp(A)}v=v$. Then
    \begin{align}
      \re\langle v,(\log A)v\rangle
      &\leq\log\!\left(\re\langle v,Av\rangle\right).
      \label{eq:mpdo_support_log_jensen}
    \end{align}
    Here $\log A$ is defined by the continuous functional calculus with
    $\log 0=0$. Compare the finite-dimensional vector-state Jensen argument
    in \cite[Chapter~V]{Bhatia1997Matrix}.
\end{theorem}

\begin{proof}\leanok
    \uses{def:hermitian_support_projection}
    Diagonalize $A=U\diag(\mu_i)U^\dagger$, set $w=U^\dagger v$, and put
    $p_i=|w_i|^2$. The support condition implies
    $\mu_i=0\Longrightarrow w_i=0$, so every spectral weight at zero vanishes.
    Thus $p_i\geq0$ for every $i\in S=\{i:\mu_i>0\}$ and
    $\sum_{i\in S}p_i=1$. The spectral formulas give
    \begin{align}
      \re\langle v,(\log A)v\rangle
      &=\sum_{i\in S}p_i\log\mu_i, \notag\\
      \re\langle v,Av\rangle
      &=\sum_{i\in S}p_i\mu_i.
      \notag
    \end{align}
    Scalar concavity of the logarithm on $(0,\infty)$ now gives
    (\ref{eq:mpdo_support_log_jensen}). No concavity assertion at zero and no
    positive-definiteness of $A$ are used.
\end{proof}

\begin{definition}[Order-two sandwiched trace functional]
    \label{def:sandwiched_renyi_two_trace}
    \lean{TNLean.sandwichedRenyiTwoTrace}
    \leanok
    For square matrices $\rho$ and $\omega$ of the same size, define
    \begin{align}
      Q_2(\rho,\omega)
      &=\re\tr\left[
        \left(\omega^{-1/4}\rho\,\omega^{-1/4}\right)^2
      \right].
      \label{eq:mpdo_sandwiched_two_trace}
    \end{align}
    The negative power is defined by functional calculus and vanishes on the
    kernel of $\omega$. On trace-one positive semidefinite matrices satisfying
    $\ker\omega\subseteq\ker\rho$, its logarithm is the order-two sandwiched
    R\'enyi divergence \cite[Definition~2]{MullerLennert2013Renyi}.
\end{definition}

\begin{theorem}[Faithful cyclic form of the order-two functional]
    \label{thm:sandwiched_renyi_two_weighted}
    \lean{TNLean.sandwichedRenyiTwoTrace_eq_weighted}
    \leanok
    \uses{def:sandwiched_renyi_two_trace}
    If $\omega\succ0$, then, for every square matrix $\rho$ of the same size,
    \begin{align}
      Q_2(\rho,\omega)
      &=\re\tr\left(\rho\,\omega^{-1/2}
        \rho\,\omega^{-1/2}\right).
      \label{eq:mpdo_sandwiched_two_weighted}
    \end{align}
\end{theorem}

\begin{proof}\leanok
    Combine the two adjacent inverse quarter-powers by
    $\omega^{-1/4}\omega^{-1/4}=\omega^{-1/2}$ and rotate the factors under
    the trace. No Hermiticity assumption on $\rho$ is needed.
\end{proof}

\begin{theorem}[Relative entropy bounded by the faithful order-two functional]
    \label{thm:relative_entropy_q2_posdef}
    \lean{TNLean.quantumRelativeEntropy_le_log_sandwichedRenyiTwoTrace_posDef}
    \leanok
    \uses{def:quantum_relative_entropy,
      def:sandwiched_renyi_two_trace}
    Let $\rho,\omega\in\MN{D}$ satisfy $\rho\succeq0$, $\omega\succ0$, and
    $\tr\rho=1$. Then
    \begin{align}
      D(\rho\Vert\omega)&\leq\log Q_2(\rho,\omega).
      \label{eq:mpdo_relative_q2_posdef}
    \end{align}
    No normalization of $\omega$ is required.
\end{theorem}

\begin{proof}\leanok
    \uses{thm:cfc_transpose_hermitian,
      thm:mpdo_support_log_kronecker,
      thm:log_sqrt_possemidef,
      thm:log_rpow_posdef,
      thm:support_log_vector_jensen,
      thm:sandwiched_renyi_two_weighted}
    Set $R=\sqrt\rho$, $T=\omega^{-1/2}$,
    $\Delta=R^T\otimes T$, and $v=\operatorname{vec}(R)$. The vectorization
    is column-stacking, with factor order
    \begin{align}
      (B\otimes A)\operatorname{vec}(X)
      &=\operatorname{vec}(AXB^T).
      \label{eq:mpdo_column_vec_kronecker}
    \end{align}
    Hence $\Delta v=\operatorname{vec}(T\rho)$, and cyclicity of the trace
    identifies its first moment with
    \begin{align}
      m=\re\langle v,\Delta v\rangle
      &=\langle R,RTR\rangle_{\mathrm{HS}}.
      \label{eq:mpdo_relative_modular_half_moment}
    \end{align}
    The support projection of $\Delta$ fixes $v$:
    \begin{align}
      P_{\supp(\Delta)}v
      &=(P_{\supp(R^T)}\otimes P_{\supp(T)})\operatorname{vec}(R)
       =\operatorname{vec}(R)=v.
      \label{eq:mpdo_relative_modular_support_fixing}
    \end{align}
    Indeed, $T$ is positive definite, so $P_{\supp(T)}=1$, while the support
    projection of $R^T$ fixes $R^T$ and hence
    $R P_{\supp(R^T)}^T=R$. The column-stacking identity
    (\ref{eq:mpdo_column_vec_kronecker}) proves the displayed equality. This
    support condition is essential when $\rho$ is singular: it removes the
    zero spectral weights before Jensen's inequality is applied.

    Theorems~\ref{thm:mpdo_support_log_kronecker},
    \ref{thm:log_sqrt_possemidef}, \ref{thm:log_rpow_posdef}, and
    \ref{thm:cfc_transpose_hermitian} give
    \begin{align}
      \re\langle v,(\log\Delta)v\rangle
      &=\frac12 D(\rho\Vert\omega).
      \label{eq:mpdo_relative_modular_log_moment}
    \end{align}
    Theorem~\ref{thm:support_log_vector_jensen} therefore yields
    $\frac12D(\rho\Vert\omega)\leq\log m$. Since
    $\|R\|_{\mathrm{HS}}^2=\tr\rho=1$, Hilbert--Schmidt
    Cauchy--Schwarz gives
    \begin{align}
      m^2
      &\leq\|R\|_{\mathrm{HS}}^2\|RTR\|_{\mathrm{HS}}^2
       =Q_2(\rho,\omega).
      \label{eq:mpdo_relative_modular_cauchy_schwarz}
    \end{align}
    Positivity of $m$ now gives
    $D(\rho\Vert\omega)\leq\log(m^2)\leq\log Q_2(\rho,\omega)$.

    The auxiliary divergence and its vector-state Jensen argument are adapted
    from \cite[arXiv:1306.3142v4, Definition~5 and Lemma~19, used in the proof
      of Theorem~7]{MullerLennert2013Renyi}.  Only the direct order-two endpoint
    is used here; the final Cauchy--Schwarz estimate is not an application of
    Beigi's interpolation theorem.
\end{proof}

\begin{theorem}[Order-two sandwiched functional under support compression]
    \label{thm:sandwiched_renyi_two_support_compression}
    \lean{TNLean.sandwichedRenyiTwoTrace_support_compression}
    \leanok
    \uses{def:sandwiched_renyi_two_trace}
    Let $\rho,\omega\in\MN{D}$ be positive semidefinite with
    $\ker\omega\subseteq\ker\rho$. Let $V:\C^k\to\C^D$ satisfy
    (\ref{eq:mpdo_reference_support_isometry}). Then
    \begin{align}
      Q_2(\rho,\omega)
      &=Q_2(V^\dagger\rho V,V^\dagger\omega V).
      \label{eq:mpdo_sandwiched_two_support_compression}
    \end{align}
\end{theorem}

\begin{proof}\leanok
    \uses{def:isometric_nonunital_conjugation,
      thm:cfc_rectangular_isometry_zero,
      thm:isometric_embedding_trace,
      thm:support_compression_reconstruction}
    The function $x\mapsto x^{-1/4}$ is taken to be zero at $x=0$, so
    Theorem~\ref{thm:cfc_rectangular_isometry_zero} applies. Reconstruct
    $\rho$ and $\omega$, combine the three expanded factors by
    multiplicativity of $\iota_V$, and use trace preservation.
\end{proof}

\begin{theorem}[Reduction of the singular order-two comparison to the faithful case]
    \label{thm:relative_entropy_q2_faithful_reduction}
    \lean{TNLean.quantumRelativeEntropy_le_log_sandwichedRenyiTwoTrace_of_faithful}
    \leanok
    \uses{def:quantum_relative_entropy,
      def:sandwiched_renyi_two_trace}
    Let $\rho,\omega\in\MN{D}$ be positive semidefinite matrices of trace one
    with $\ker\omega\subseteq\ker\rho$. Assume that for every dimension and
    every pair of trace-one matrices $A\succeq0$ and $B\succ0$,
    \begin{align}
      D(A\Vert B)&\leq\log Q_2(A,B).
      \label{eq:mpdo_faithful_relative_q2_hypothesis}
    \end{align}
    Then
    \begin{align}
      D(\rho\Vert\omega)&\leq\log Q_2(\rho,\omega).
      \label{eq:mpdo_singular_relative_q2_comparison}
    \end{align}
    Thus this theorem reduces the singular-reference case to the faithful
    comparison. Theorem~\ref{thm:relative_entropy_q2_posdef} supplies the
    hypothesis in~(\ref{eq:mpdo_faithful_relative_q2_hypothesis}).
\end{theorem}

\begin{proof}\leanok
    \uses{thm:compression_on_support_posDef,
      thm:isometric_embedding_trace,
      thm:support_compression_reconstruction,
      thm:relative_entropy_support_compression,
      thm:sandwiched_renyi_two_support_compression}
    Choose an isometric inclusion of the support of $\omega$. Its compression
    $\omega_V=V^\dagger\omega V$ is positive definite by
    Theorem~\ref{thm:compression_on_support_posDef}, while
    $\rho_V=V^\dagger\rho V$ remains positive semidefinite. Reconstruction
    and (\ref{eq:mpdo_isometric_trace}) show that both compressed matrices
    have trace one. Apply (\ref{eq:mpdo_faithful_relative_q2_hypothesis}) to
    $(\rho_V,\omega_V)$, then use
    (\ref{eq:mpdo_relative_entropy_support_compression}) and
    (\ref{eq:mpdo_sandwiched_two_support_compression}).
\end{proof}

\begin{theorem}[Relative entropy bounded by the support-restricted order-two functional]
    \label{thm:relative_entropy_q2_support}
    \lean{TNLean.quantumRelativeEntropy_le_log_sandwichedRenyiTwoTrace}
    \leanok
    \uses{def:quantum_relative_entropy,
      def:sandwiched_renyi_two_trace}
    Let $\rho,\omega\in\MN{D}$ be positive semidefinite matrices of trace one.
    If $\ker\omega\subseteq\ker\rho$, then
    \begin{align}
      D(\rho\Vert\omega)&\leq\log Q_2(\rho,\omega).
      \label{eq:mpdo_relative_q2_support}
    \end{align}
\end{theorem}

\begin{proof}\leanok
    \uses{thm:relative_entropy_q2_posdef,
      thm:relative_entropy_q2_faithful_reduction}
    Apply Theorem~\ref{thm:relative_entropy_q2_faithful_reduction} with the
    faithful comparison from Theorem~\ref{thm:relative_entropy_q2_posdef}.
    Compression to the support of $\omega$ preserves both entropy functionals
    and both trace normalizations, and makes the reference positive definite.
\end{proof}

\section{Rectangular Choi matrices and a range-dimension estimate}

The following construction is the unnormalized, different-dimension form of
the Choi correspondence \cite[Proposition~2.1]{Wolf2012Quantum}. The resulting
estimate bounds the order-two Choi functional of a Hilbert--Schmidt contraction
by the dimension of its range. It is used below to bound mutual information in
terms of operator-Schmidt rank.

\begin{definition}[Matrix of a linear map in matrix-unit bases]
    \label{def:linear_map_matrix}
    \lean{Matrix.linearMapMatrix}
    \leanok
    Let $\mathcal L:\MN{d}\to\MN{e}$ be complex-linear. Its coefficient
    matrix $C(\mathcal L)\in M_{e^2\times d^2}(\C)$ in the matrix-unit bases is
    defined by
    \begin{align}
        C(\mathcal L)_{(a,b),(i,j)}
        &=(\mathcal L(E_{ij}))_{ab}.
        \label{eq:rectangular_choi_coefficient_matrix}
    \end{align}
\end{definition}

\begin{definition}[Unnormalized rectangular Choi matrix]
    \label{def:rectangular_choi}
    \lean{Matrix.rectangularChoi}
    \leanok
    \uses{def:linear_map_matrix}
    For a complex-linear map $\mathcal L:\MN{d}\to\MN{e}$, its unnormalized
    rectangular Choi matrix $J(\mathcal L)\in\MN{de}$ is the reshaping
    \begin{align}
        J(\mathcal L)_{(i,a),(j,b)}
        &=C(\mathcal L)_{(a,b),(i,j)}
         =(\mathcal L(E_{ij}))_{ab}.
        \label{eq:rectangular_choi_entries}
    \end{align}
    When $d=e$, this is $d$ times the normalized Choi matrix convention in
    \cite[Proposition~2.1]{Wolf2012Quantum}.
\end{definition}

\begin{theorem}[Frobenius norm under rectangular Choi reshaping]
    \label{thm:rectangular_choi_parseval}
    \lean{Matrix.rectangularChoi_frobenius_parseval}
    \leanok
    \uses{def:linear_map_matrix, def:rectangular_choi}
    For every complex-linear map $\mathcal L:\MN{d}\to\MN{e}$,
    \begin{align}
        \re\tr\!\left(J(\mathcal L)^\dagger J(\mathcal L)\right)
        &=
        \re\tr\!\left(C(\mathcal L)^\dagger C(\mathcal L)\right).
        \label{eq:rectangular_choi_parseval}
    \end{align}
\end{theorem}

\begin{proof}\leanok
    \uses{def:rectangular_choi}
    The two traces are the sums of the squared absolute values of their
    respective matrix entries. The bijection
    \begin{align}
        ((i,a),(j,b))\longmapsto((i,j),(a,b))
        \notag
    \end{align}
    identifies the two sums by~(\ref{eq:rectangular_choi_entries}).
\end{proof}

\begin{theorem}[Matrix-unit Parseval identity for a rectangular Choi matrix]
    \label{thm:rectangular_choi_parseval_matrix_units}
    \lean{Matrix.rectangularChoi_frobenius_parseval_matrix_units}
    \leanok
    \uses{def:rectangular_choi}
    For every complex-linear map $\mathcal L:\MN{d}\to\MN{e}$,
    \begin{align}
        \re\tr\!\left(J(\mathcal L)^\dagger J(\mathcal L)\right)
        &=
        \sum_{i,j=0}^{d-1}\sum_{a,b=0}^{e-1}
        \left|(\mathcal L(E_{ij}))_{ab}\right|^2.
        \label{eq:rectangular_choi_parseval_matrix_units}
    \end{align}
\end{theorem}

\begin{proof}\leanok
    \uses{def:linear_map_matrix, thm:rectangular_choi_parseval}
    Apply Theorem~\ref{thm:rectangular_choi_parseval} and substitute
    (\ref{eq:rectangular_choi_coefficient_matrix}) into the entrywise
    Frobenius-norm formula for $C(\mathcal L)$.
\end{proof}

\begin{definition}[Hilbert--Schmidt contraction]
    \label{def:rectangular_hilbert_schmidt_contraction}
    \lean{Matrix.IsHilbertSchmidtContraction}
    \leanok
    A complex-linear map $\mathcal L:\MN{d}\to\MN{e}$ is a
    \emph{Hilbert--Schmidt contraction} if, for every $X\in\MN{d}$,
    \begin{align}
        \re\tr\!\left((\mathcal L(X))^\dagger\mathcal L(X)\right)
        &\leq \re\tr(X^\dagger X).
        \label{eq:rectangular_hilbert_schmidt_contraction}
    \end{align}
\end{definition}

\begin{theorem}[Rectangular Choi norm bounded by the range dimension]
    \label{thm:rectangular_choi_rank_bound}
    \lean{Matrix.rectangularChoi_frobenius_sq_le_finrank_range}
    \leanok
    \uses{def:rectangular_choi,
      def:rectangular_hilbert_schmidt_contraction}
    If $\mathcal L:\MN{d}\to\MN{e}$ is a Hilbert--Schmidt contraction, then
    \begin{align}
        \left\|J(\mathcal L)\right\|_F^2
        &\leq\dim_{\C}\range\mathcal L.
        \label{eq:rectangular_choi_rank_bound}
    \end{align}
\end{theorem}

\begin{proof}\leanok
    \uses{def:linear_map_matrix, thm:rectangular_choi_parseval}
    Set $G=C(\mathcal L)^\dagger C(\mathcal L)$. The contraction inequality
    applied in the matrix-unit bases shows that every eigenvalue of the
    positive semidefinite matrix $G$ is at most one. Hence
    \begin{align}
        \tr G
        &\leq\#\{k:\lambda_k(G)\neq0\}
         =\rank C(\mathcal L)
         =\dim_{\C}\range\mathcal L.
        \notag
    \end{align}
    Theorem~\ref{thm:rectangular_choi_parseval} identifies the left-hand side
    with $\|J(\mathcal L)\|_F^2$.
\end{proof}

\begin{theorem}[Full-support eigenbasis whitening as a Choi matrix]
    \label{thm:full_support_eigenbasis_whitening_choi}
    \lean{Matrix.rectangularChoi_singleKrausMap_comp_comp,
      Matrix.supportedMarginalWhitenedState_eq_rectangularChoi,
      Matrix.finrank_range_weightedHilbertSchmidtMap}
    \leanok
    \uses{thm:supported_marginal_channel,
      def:rectangular_choi,
      def:weighted_hilbert_schmidt_map}
    Let $\rho_{AB}\geq0$, and choose an eigenbasis of its faithful first
    marginal
    $\sigma=\diag(p_1,\ldots,p_{d_A})>0$.  Suppose also that the second
    marginal $\tau$ is positive definite.  For the supported-marginal channel
    $\Phi_\rho$, set
    \begin{align}
      L(X)&=\tau^{-1/4}
        \Phi_\rho(\sigma^{1/4}X\sigma^{1/4})\tau^{-1/4},\\
      W&=(\sigma^T\otimes\tau)^{-1/4}\rho_{AB}
        (\sigma^T\otimes\tau)^{-1/4}.
        \notag
    \end{align}
    Then
    \begin{align}
      W&=J(L),
      &\dim_{\C}\range L&=\dim_{\C}\range\Phi_\rho.
      \notag
    \end{align}
    Here $\sigma^T=\sigma$ because the chosen basis diagonalizes $\sigma$.
\end{theorem}

\begin{proof}\leanok
    \uses{thm:supported_marginal_channel}
    For arbitrary congruence matrices $A$ and $B$, the matrix-unit convention
    for the Choi matrix gives
    \begin{align}
      J(\Gamma_B\circ\Phi\circ\Gamma_A)
      &=\Gamma_{A^T\otimes B}(J(\Phi)).
      \notag
    \end{align}
    Apply this identity to the inverse-square-root scaling in
    $\Phi_\rho$ and to the two fourth-root factors defining $L$.
    Entrywise diagonal functional calculus gives
    $\sigma^{1/4}=\diag(p_i^{1/4})$ and
    $\diag(p_i^{-1/2})\sigma^{1/4}=\sigma^{-1/4}$, proving the first
    assertion.  The two modular congruences are invertible linear maps, so
    pre- and post-composition by them preserve the range dimension.
\end{proof}

\begin{theorem}[Full-support eigenbasis whitened Choi estimate]
    \label{thm:full_support_eigenbasis_whitened_choi_bound}
    \lean{Matrix.supportedMarginalWhitenedState_posSemidef,
      Matrix.supportedMarginalWhitenedState_frobenius_sq_le_operatorSchmidtRank,
      Matrix.supportedMarginalWhitenedState_trace_sq_re_le_operatorSchmidtRank}
    \leanok
    \uses{thm:full_support_eigenbasis_whitening_choi,
      thm:rectangular_choi_rank_bound,
      thm:weighted_hilbert_schmidt_contraction_full_support}
    Under the hypotheses of
    Theorem~\ref{thm:full_support_eigenbasis_whitening_choi},
    the matrix
    \begin{align}
      W=(\sigma^T\otimes\tau)^{-1/4}\rho_{AB}
        (\sigma^T\otimes\tau)^{-1/4}
        \notag
    \end{align}
    is positive semidefinite, and
    \begin{align}
      \tr(W^2)=\|W\|_F^2
      &\leq \operatorname{OSR}(\rho_{AB}).
      \notag
    \end{align}
    This is the faithful-marginal-support form of the order-two estimate
    obtained from \cite[Theorem~6, Equation~(18)]{Beigi2013Sandwiched}.
    % The passage from singular ambient marginals to their supports remains
    % recorded in docs/paper-gaps/cpgsv17_mpdo_mutual_information_bound.tex.
\end{theorem}

\begin{proof}\leanok
    \uses{thm:full_support_eigenbasis_whitening_choi,
      thm:rectangular_choi_rank_bound,
      thm:weighted_hilbert_schmidt_contraction_full_support}
    The matrix $W$ is a congruence of $\rho_{AB}\geq0$, so it is positive
    semidefinite and Hermitian; consequently
    $\tr(W^2)=\|W\|_F^2$.
    The weighted map $L$ is a Hilbert--Schmidt contraction.  Its Choi matrix
    is the whitened state by the preceding theorem, while invertibility of
    the modular congruences and the supported-marginal correspondence give
    \begin{align}
      \dim_{\C}\range L
      &=\dim_{\C}\range\Phi_\rho
       =\operatorname{OSR}(\rho_{AB}).
      \notag
    \end{align}
    The rectangular Choi range-dimension estimate now proves the claim.
\end{proof}

\begin{theorem}[Real powers under unitary conjugation]
    \label{thm:rpow_conj_unitary}
    \lean{Matrix.rpow_conj_unitary}
    \leanok
    \uses{lem:cfc_conj_unitary}
    Let $A\succeq0$, let $s\in\R$, and let $U$ be unitary. Then
    \begin{align}
      (UAU^\dagger)^s&=UA^sU^\dagger.
      \notag
    \end{align}
\end{theorem}

\begin{proof}\leanok
    \uses{lem:cfc_conj_unitary}
    Apply Lemma~\ref{lem:cfc_conj_unitary} to the function $x\mapsto x^s$.
\end{proof}

\begin{theorem}[Real powers of positive tensor products]
    \label{thm:rpow_kronecker_possemidef}
    \lean{Matrix.PosSemidef.rpow_kronecker}
    \leanok
    \uses{thm:multiplicative_functional_calculus_tensor_product}
    Let $A\succeq0$ and $B\succeq0$. For every $s\in\R$,
    \begin{align}
      (A\otimes B)^s&=A^s\otimes B^s.
      \notag
    \end{align}
\end{theorem}

\begin{proof}\leanok
    \uses{thm:multiplicative_functional_calculus_tensor_product}
    Apply
    Theorem~\ref{thm:multiplicative_functional_calculus_tensor_product} to
    $x\mapsto x^s$, using $(xy)^s=x^sy^s$ for $x,y\geq0$.
\end{proof}

\begin{theorem}[Inverse quarter-power from iterated square roots]
    \label{thm:rpow_neg_quarter_inv_sqrt_sqrt}
    \lean{Matrix.PosDef.rpow_neg_quarter_eq_inv_sqrt_sqrt}
    \leanok
    If $A\succ0$, then
    \begin{align}
      A^{-1/4}&=\left(\sqrt{\sqrt A}\right)^{-1}.
      \notag
    \end{align}
\end{theorem}

\begin{proof}\leanok
    The continuous functional calculus gives
    $\sqrt{\sqrt A}=A^{1/4}$. Positive definiteness makes this matrix
    invertible, and inversion changes the exponent from $1/4$ to $-1/4$.
\end{proof}

\begin{theorem}[Unitary invariance of the order-two functional]
    \label{thm:sandwiched_renyi_two_unitary_invariance}
    \lean{TNLean.sandwichedRenyiTwoTrace_conj_unitary}
    \leanok
    \uses{def:sandwiched_renyi_two_trace, thm:rpow_conj_unitary}
    Let $\omega\succeq0$, let $\rho$ be a square matrix of the same size, and
    let $U$ be unitary. Then
    \begin{align}
      Q_2(U\rho U^\dagger,U\omega U^\dagger)
      &=Q_2(\rho,\omega).
      \notag
    \end{align}
\end{theorem}

\begin{proof}\leanok
    \uses{thm:rpow_conj_unitary}
    Theorem~\ref{thm:rpow_conj_unitary} gives
    $(U\omega U^\dagger)^{-1/4}=U\omega^{-1/4}U^\dagger$.
    Substitute this identity into both sandwiching factors and use
    $U^\dagger U=1$ and cyclicity of the trace.
\end{proof}

\begin{theorem}[Faithful marginal whitening bound]
    \label{thm:sandwiched_renyi_two_osr_faithful}
    \lean{TNLean.sandwichedRenyiTwoTrace_partialTraces_kronecker_le_operatorSchmidtRank_of_posDef}
    \leanok
    \uses{def:sandwiched_renyi_two_trace,
      def:bipartite_operator_schmidt_rank,
      lem:partial_trace_product_isometry,
      thm:operator_schmidt_rank_local_isometry,
      thm:full_support_eigenbasis_whitened_choi_bound,
      thm:rpow_kronecker_possemidef,
      thm:rpow_neg_quarter_inv_sqrt_sqrt,
      thm:sandwiched_renyi_two_unitary_invariance}
    Let $\rho_{AB}\succeq0$ have positive definite marginals $\rho_A$ and
    $\rho_B$. Then
    \begin{align}
      Q_2(\rho_{AB},\rho_A\otimes\rho_B)
      &\leq\operatorname{OSR}(\rho_{AB}).
      \label{eq:mpdo_faithful_q2_osr}
    \end{align}
    No trace normalization is required.
\end{theorem}

\begin{proof}\leanok
    \uses{lem:partial_trace_product_isometry,
      thm:operator_schmidt_rank_local_isometry,
      thm:full_support_eigenbasis_whitened_choi_bound,
      thm:rpow_kronecker_possemidef,
      thm:rpow_neg_quarter_inv_sqrt_sqrt,
      thm:sandwiched_renyi_two_unitary_invariance}
    Choose a unitary $U$ that diagonalizes the first marginal and conjugate
    $\rho_{AB}$ by $U^\dagger\otimes1$. The two marginals become
    $\diag(p_i)$ and $\rho_B$, while positivity and operator-Schmidt rank are
    unchanged. The transpose in the Choi whitening convention becomes
    invisible only in this eigenbasis, since $\diag(p_i)^T=\diag(p_i)$.

    Theorems~\ref{thm:rpow_kronecker_possemidef} and
    \ref{thm:rpow_neg_quarter_inv_sqrt_sqrt} identify the negative
    quarter-power of the product marginal with the two whitening factors in
    Theorem~\ref{thm:full_support_eigenbasis_whitened_choi_bound}. That theorem
    bounds the squared whitened trace by the ordinary operator-Schmidt rank.
    Finally, Theorem~\ref{thm:sandwiched_renyi_two_unitary_invariance} and the
    invariance of operator-Schmidt rank under local unitaries return
    to the original basis.
\end{proof}

\begin{theorem}[Marginal-support whitening bound]
    \label{thm:sandwiched_renyi_two_osr_support}
    \lean{TNLean.sandwichedRenyiTwoTrace_partialTraces_kronecker_le_operatorSchmidtRank}
    \leanok
    \uses{thm:sandwiched_renyi_two_osr_faithful,
      thm:sandwiched_renyi_two_support_compression,
      thm:operator_schmidt_rank_marginal_support_compression,
      thm:faithful_marginals_support_compression}
    Let $\rho_{AB}\succeq0$ be any finite-dimensional bipartite operator. Then
    \begin{align}
      Q_2(\rho_{AB},\rho_A\otimes\rho_B)
      &\leq\operatorname{OSR}(\rho_{AB}).
      \label{eq:mpdo_support_q2_osr}
    \end{align}
    No trace normalization or positive-dimension hypothesis is required. The
    statement includes the cases in which either marginal support has
    dimension zero.
\end{theorem}

\begin{proof}\leanok
    \uses{thm:sandwiched_renyi_two_osr_faithful,
      thm:sandwiched_renyi_two_support_compression,
      thm:operator_schmidt_rank_marginal_support_compression,
      thm:faithful_marginals_support_compression}
    Compress both factors simultaneously to the supports of $\rho_A$ and
    $\rho_B$. The compressed marginals are positive definite by
    Theorem~\ref{thm:faithful_marginals_support_compression}, even when a
    support coordinate space has dimension zero. The order-two functional is
    unchanged by
    Theorem~\ref{thm:sandwiched_renyi_two_support_compression}, and the
    ordinary operator-Schmidt rank is unchanged by
    Theorem~\ref{thm:operator_schmidt_rank_marginal_support_compression}.
    Apply Theorem~\ref{thm:sandwiched_renyi_two_osr_faithful} to the compressed
    operator.
\end{proof}

The sandwiched R\'enyi comparison also uses the following unnormalized
trace term.  For trace-one $\rho$ and $\alpha>1$, it is the quantity inside the
logarithm of the sandwiched R\'enyi divergence
\cite[Equation~(3)]{Beigi2013Sandwiched}
\cite[Definition~2]{MullerLennert2013Renyi} whenever
$\rho,\omega\geq0$ and $\ker\omega\subseteq\ker\rho$.  This includes singular
$\omega$, with negative powers interpreted as the generalized inverse on its
support.  In this regime, finiteness of the divergence requires the support
inclusion.  The present application uses $1<\alpha\leq2$; the order-one case is
the Umegaki endpoint.  The statements below concern only this algebraic
expression.  They are not assertions of
\cite[Proposition~4.5]{Cirac2016MPDO_arXiv}, and they do not establish
continuity or monotonicity in the order, interpolation between the endpoints,
differentiability at order one, a limit to Umegaki relative entropy, a
logarithmic divergence, or a comparison with relative entropy.

\begin{definition}[Total sandwiched R\'enyi trace]
    \label{def:sandwiched_renyi_trace}
    \lean{TNLean.sandwichedRenyiTrace}
    \leanok
    For square complex matrices $\rho$ and $\omega$ and every $\alpha\in\R$,
    set
    \begin{align}
        \widetilde Q_\alpha(\rho\Vert\omega)
        &=\re\tr\!\left[
          \left(\omega^{(1-\alpha)/(2\alpha)}
          \rho\,\omega^{(1-\alpha)/(2\alpha)}\right)^\alpha
          \right].
        \label{eq:sandwiched_renyi_trace}
    \end{align}
    Real powers are defined by continuous functional calculus, with negative
    powers equal to zero on the zero eigenspace.  Thus
    $\widetilde Q_\alpha$ is total even when $\alpha=0$ or $\omega$ is singular.
    For $\alpha>1$, on positive semidefinite inputs satisfying
    $\ker\omega\subseteq\ker\rho$, this convention gives the finite
    sandwiched R\'enyi trace term.  In this regime, if the support inclusion
    fails, the totalized value remains finite but is not the divergence trace
    term.
\end{definition}

\begin{theorem}[Nonnegativity on the faithful domain]
    \label{thm:sandwiched_renyi_trace_nonneg}
    \lean{TNLean.sandwichedRenyiTrace_nonneg}
    \leanok
    \uses{def:sandwiched_renyi_trace}
    If $\rho\geq0$ and $\omega>0$, then, for every $\alpha\in\R$,
    \begin{align}
        0&\leq\widetilde Q_\alpha(\rho\Vert\omega).
        \notag
    \end{align}
\end{theorem}

\begin{proof}\leanok
    \uses{def:sandwiched_renyi_trace}
    Put $q=\omega^{(1-\alpha)/(2\alpha)}$.  Continuous functional calculus
    gives $q\geq0$, hence $q\rho q\geq0$ and $(q\rho q)^\alpha\geq0$.
    The trace of the last matrix is real and nonnegative.
\end{proof}

\begin{theorem}[Order-one sandwiched trace]
    \label{thm:sandwiched_renyi_trace_one}
    \lean{TNLean.sandwichedRenyiTrace_one}
    \leanok
    \uses{def:sandwiched_renyi_trace}
    If $\rho\geq0$ and $\omega>0$, then
    \begin{align}
        \widetilde Q_1(\rho\Vert\omega)&=\re\tr(\rho).
        \notag
    \end{align}
\end{theorem}

\begin{proof}\leanok
    \uses{def:sandwiched_renyi_trace}
    At $\alpha=1$ the two powers of $\omega$ have exponent zero, while the
    remaining real power of $\rho$ has exponent one.  Substitution into
    (\ref{eq:sandwiched_renyi_trace}) gives the identity.
\end{proof}

\begin{theorem}[Order-two sandwiched trace]
    \label{thm:sandwiched_renyi_trace_two}
    \lean{TNLean.sandwichedRenyiTrace_two}
    \leanok
    \uses{def:sandwiched_renyi_trace,
      def:sandwiched_renyi_two_trace}
    If $\rho\geq0$ and $\omega>0$, then
    \begin{align}
        \widetilde Q_2(\rho\Vert\omega)&=Q_2(\rho,\omega).
        \notag
    \end{align}
\end{theorem}

\begin{proof}\leanok
    \uses{def:sandwiched_renyi_trace,
      def:sandwiched_renyi_two_trace}
    At $\alpha=2$ the sandwiching exponent is $-1/4$.
    Since $\omega^{-1/4}\rho\,\omega^{-1/4}$ is positive semidefinite, its
    continuous-functional-calculus square is its ordinary matrix square.
\end{proof}

\begin{theorem}[Mutual information from operator-Schmidt rank]
    \label{thm:mutual_information_le_log_operator_schmidt_rank}
    \uses{def:bipartite_operator_schmidt_rank}
    For every finite-dimensional bipartite density operator $\rho_{AB}$,
    \begin{align}
        I(A:B)_\rho
        \leq \log\operatorname{OSR}(\rho_{AB}).
        \notag
    \end{align}
    % Formalization is tracked by issue #5212.
\end{theorem}
\begin{proof}
    \uses{thm:relative_entropy_q2_support,
      thm:sandwiched_renyi_two_osr_support}
    Set $\omega=\rho_A\otimes\rho_B$. Positivity gives
    $\ker\omega\subseteq\ker\rho_{AB}$, and the two matrices have trace one.
    Therefore
    \begin{align}
      I(A:B)_\rho
      &=D(\rho_{AB}\Vert\omega) \\
      &\leq\log Q_2(\rho_{AB},\omega) \\
      &\leq\log\operatorname{OSR}(\rho_{AB}).
      \notag
    \end{align}
    The first inequality is
    Theorem~\ref{thm:relative_entropy_q2_support}; the second follows from
    Theorem~\ref{thm:sandwiched_renyi_two_osr_support} and monotonicity of the
    logarithm.
\end{proof}

\begin{remark}[Boundedness of the mutual information]
    The proof of Proposition~4.5 of \cite{Cirac2016MPDO_arXiv} invokes the
    uniform estimate $I_L\leq4\log D$ for every MPDO. Its cited argument,
    however, treats mixed tensor-network states obtained from local completely
    positive maps \cite{Wolf2008AreaLaws}. After purifying those maps, data
    processing and the pure-state boundary estimate give
    \begin{align}
        I(A:B)_{\rho}
        &\leq I(A:B)_{\ketbra{\Psi}{\Psi}}
         =2S(A)_{\ketbra{\Psi}{\Psi}}
         \leq2\lvert\partial A\rvert\log D.
        \notag
    \end{align}
    A periodic one-dimensional interval has two boundary bonds, and hence the
    last expression is $4\log D$. A general positive MPO need not admit this
    local purification.

    By Theorem~\ref{thm:operator_schmidt_coefficient_space}, the two virtual
    bonds give an operator-Schmidt decomposition with at most $D^2$ terms, but
    they do not bound the ordinary rank of either marginal.
    For example, a bond-one tensor can generate $q^{\otimes N}$ for a full-rank
    one-site density matrix $q$. Thus its $L$-site marginal has rank
    $\rank(q)^L$, although its mutual information vanishes.

    Theorem~\ref{thm:mutual_information_le_log_operator_schmidt_rank} instead
    gives the stronger estimate
    \begin{align}
        I_L^{(N)}\leq\log D^2=2\log D
        \notag
    \end{align}
    whenever the finite-chain operator is positive and has positive trace.
    Thus no classical or genuinely quantum counterexample exists.
    The diagonal argument gives an independent special-case proof. Indeed,
    Theorem~\ref{thm:mpo_diagonal_finite_chain_classical_mutual_information}
    bounds by $2\log D$ the classical mutual information of the normalized
    diagonal coefficients whenever the finite-chain operator is diagonal and
    positive semidefinite and has positive total mass.

    The same example as
    Theorem~\ref{thm:parity_mpdo_no_thermodynamic_limit} makes the iterated
    mutual-information limit false; see
    \path{docs/paper-gaps/cpgsv17_mpdo_mutual_information_bound.tex}. The
    channel-image implication is
    Theorem~\ref{thm:local_channel_image_mutual_info_bound}, and the pure-state
    instance is Proposition~\ref{prop:pure_mutual_info_bound} below. The
    locally purified case is Theorem~\ref{thm:lpdo_mutual_info_bound}.
\end{remark}

\begin{definition}[Block entropy and area-law saturation for pure states]
    \label{def:pure_sal}
    \lean{MPSTensor.pureBlockEntropy, MPSTensor.IsSAL}
    \leanok
    \uses{def:mps_tensor, def:von_neumann_entropy}
    For an MPS tensor $A$, a system size $N$, and a block length $L\le N$, let
    \begin{align}
        \sigma^{(N)}(A)
        &=
        \frac{\ketbra{V^{(N)}(A)}{V^{(N)}(A)}}{\|V^{(N)}(A)\|^2},
        \notag\\
        \rho_L^{(N)}(A)
        &=\tr_{N\setminus L}\![\sigma^{(N)}(A)]
        \notag
    \end{align}
    denote the normalized state and the reduced state of the first $L$ spins.
    The $L$-block entropy is
    $S_L^{(N)}(A)=S(\rho_L^{(N)}(A))$.
    The tensor saturates the area law if
    $S_L^{(N)}(A)=S_{L+1}^{(N)}(A)$ for all
    $1\le L<\lfloor N/2\rfloor$ and all $N$.
    This is \cite[Definition~3.13, line~600]{Cirac2016MPDO_arXiv}.
\end{definition}

\begin{definition}[Purification tensor]\label{def:doubled_tensor}
    \lean{doubledTensor}
    \leanok
    \uses{def:mps_tensor, def:mpo_tensor}
    For an MPS tensor $A$ on bond space $\C^D$, define the associated
    MPO tensor $M$ on bond space $\C^{D^2}$ by
    $M^{ij}=A^i\otimes\overline{A^j}$.
    This is the single-ancilla purification tensor of
    \cite{Cirac2016MPDO_arXiv}.
\end{definition}

\begin{lemma}[Pure state as a purification MPDO]\label{lem:mpo_doubled}
    \lean{mpo_doubledTensor}
    \leanok
    \uses{def:mps_tensor, def:mpo_family, def:doubled_tensor}
    For an MPS tensor $A$, the pure-state operator
    $\ketbra{V^{(N)}(A)}{V^{(N)}(A)}$ equals the MPDO generated by the tensor
    $M^{ij}=A^i\otimes\overline{A^j}$, acting on a bond space of
    dimension $D^2$:
    \begin{align}
        \rho^{(N)}(A\otimes\bar A)
        &=\ketbra{V^{(N)}(A)}{V^{(N)}(A)}.
        \notag
    \end{align}
    This is the single-ancilla case of the purification picture of
    \cite{Cirac2016MPDO_arXiv}; it lets the MPDO block-entropy
    theory apply to pure-state block entropies.
\end{lemma}
\begin{proof}\leanok
    The word evaluation of the doubled tensor is the Kronecker product of the
    word evaluation of $A$ and its conjugate,
    $M^{\sigma,\tau}=A^\sigma\otimes\overline{A^\tau}$,
    so its trace factors as
    $\tr(M^{\sigma,\tau})
      =\tr(A^\sigma) \overline{\tr(A^\tau)}
      =\overline{V^{(N)}(A)_\tau}\,V^{(N)}(A)_\sigma$,
    the corresponding matrix element of
    $\ketbra{V^{(N)}(A)}{V^{(N)}(A)}$.
\end{proof}

\begin{lemma}[Block entropy of the purification tensor]
    \label{lem:blockEntropy_doubledTensor}
    \lean{blockEntropy_doubledTensor}
    \leanok
    \uses{lem:mpo_doubled, def:block_entropy, def:pure_sal}
    For an MPS tensor $A$, system size $N$, and block length $L\le N$, the MPDO
    block entropy of the tensor $M^{ij}=A^i\otimes\overline{A^j}$ agrees
    with the pure-state block entropy:
    $S_L^{(N)}(M)=S_L^{(N)}(A)$.
\end{lemma}
\begin{proof}\leanok
    Lemma~\ref{lem:mpo_doubled} gives
    $\rho^{(N)}(M)=\ketbra{V^{(N)}(A)}{V^{(N)}(A)}$, hence
    \begin{align}
        \sigma^{(N)}(M)
        &=
        \frac{\rho^{(N)}(M)}{\tr[\rho^{(N)}(M)]}
        =
        \frac{\ketbra{V^{(N)}(A)}{V^{(N)}(A)}}{\|V^{(N)}(A)\|^2}
        =
        \sigma^{(N)}(A).
        \notag
    \end{align}
    Therefore the reduced block states agree:
    \begin{align}
        \rho_L^{(N)}(M)
        &=
        \tr_{N\setminus L}\![\sigma^{(N)}(M)]
        =
        \tr_{N\setminus L}\![\sigma^{(N)}(A)]
        =
        \rho_L^{(N)}(A).
        \notag
    \end{align}
    Thus
    $S_L^{(N)}(M)
      =S(\rho_L^{(N)}(M))
      =S(\rho_L^{(N)}(A))
      =S_L^{(N)}(A)$.
\end{proof}

\begin{lemma}[Schmidt symmetry of the block entropy]\label{lem:schmidt_symmetry}
    \lean{pureBlockEntropy_complement}
    \leanok
    \uses{def:pure_sal}
    For an MPS tensor $A$ and a block length $L\le N$, the block entropies of a
    block and its complement coincide:
    $S_L^{(N)}(A)=S_{N-L}^{(N)}(A)$.
\end{lemma}
\begin{proof}\leanok
    The reduced state of the first $L$ spins is $\rho_L\propto WW^\dagger$, where
    $W$ is the matrix of amplitudes $\bra{u\,w}V^{(N)}(A)\rangle$ with $u$ the
    first $L$ spins and $w$ the remaining $N-L$. By cyclic invariance of the
    amplitudes, the reduced state of the complement is
    $\propto\overline{W^\dagger W}$.
    Cyclic invariance of the entropy,
    $S(WW^\dagger)=S(W^\dagger W)$, together with invariance of the entropy
    under entrywise conjugation of a Hermitian matrix,
    $S(\overline{W^\dagger W})=S(W^\dagger W)$, gives
    $S_L^{(N)}(A)=S_{N-L}^{(N)}(A)$.
\end{proof}

\begin{proposition}[Pure-state area-law monotonicity]\label{prop:pure_area_monotone}
    \lean{pureBlockEntropy_monotone}
    \leanok
    \uses{def:pure_sal, lem:schmidt_symmetry, lem:blockEntropy_doubledTensor}
    Let $A$ be an MPS tensor with $V^{(N)}(A)\neq0$. For every block length $L$
    with $2L+1\le N$, covering $1\le L<\lfloor N/2\rfloor$,
    $S_L^{(N)}(A)\le S_{L+1}^{(N)}(A)$.
    This is the pure-state area law of
    \cite[Section~3, line~599]{Cirac2016MPDO_arXiv}.
\end{proposition}
\begin{proof}\leanok
    \uses{prop:mutual_info_monotone}
    Apply the block-entropy form of strong subadditivity, the inequality
    underlying Proposition~\ref{prop:mutual_info_monotone}, to the purification
    MPDO of $A$, giving
    $S_{N-L}+S_L\le S_{L+1}+S_{N-L-1}$. By the Schmidt
    symmetry $S_{N-L}=S_L$ and $S_{N-L-1}=S_{L+1}$, this reduces to
    $S_L\le S_{L+1}$.
\end{proof}

\begin{lemma}[Operator-Schmidt Gram is a transfer-map power]
    \label{lem:schmidt_gram_transfer_power}
    \lean{MPSTensor.schmidtLeft, MPSTensor.schmidtRight,
        MPSTensor.schmidtLeft_gram_apply, MPSTensor.schmidtRight_gram_apply}
    \leanok
    \uses{def:transfer_map}
    Let $A$ be an MPS tensor of bond dimension $D$. For a length-$\ell$
    configuration $\sigma$ and a length-$m$ configuration $\tau$, write the
    operator-Schmidt factors as
    \begin{align}
        (L_\ell)_{\sigma,(a_1,a_2)}
        &=(A^\sigma)_{a_1a_2},
        \notag\\
        (R_m)_{(p_1,p_2),\tau}
        &=(A^\tau)_{p_2p_1}.
        \notag
    \end{align}
    The $D^2\times D^2$ Gram matrix of the left factor equals the $\ell$-fold
    transfer map evaluated on a matrix unit:
    \begin{align}
        (L_\ell^\dagger L_\ell)_{(a_1,a_2),(b_1,b_2)}
        &=
        \sum_{\sigma}\overline{(A^\sigma)_{a_1a_2}}
            (A^\sigma)_{b_1b_2}
        \notag\\
        &=
        (\E_A^{\,\ell}(\ketbra{b_2}{a_2}))_{b_1,a_1},
        \notag
    \end{align}
    where the sum runs over length-$\ell$ configurations $\sigma$. The Gram of
    the right factor $R_m$ of an $m$-spin complement satisfies the bond-swapped
    identity
    \begin{align}
        (R_mR_m^\dagger)_{(p_1,p_2),(q_1,q_2)}
        &=
        \sum_{\tau}(A^\tau)_{p_2p_1}
            \overline{(A^\tau)_{q_2q_1}}
        \notag\\
        &=
        (\E_A^{\,m}(\ketbra{p_1}{q_1}))_{p_2,q_2},
        \notag
    \end{align}
    where the sum runs over length-$m$ configurations $\tau$.
\end{lemma}
\begin{proof}\leanok
    The left Gram entry is
    $\sum_\sigma\overline{(A^\sigma)_{a_1a_2}}(A^\sigma)_{b_1b_2}$
    by definition of $L_\ell$. Expanding
    \begin{align}
        \E_A^{\,\ell}(\ketbra{b_2}{a_2})
        &=
        \sum_\sigma
        A^\sigma\ketbra{b_2}{a_2}(A^\sigma)^\dagger
        \notag
    \end{align}
    and reading off the $(b_1,a_1)$ entry gives the same sum, since
    \begin{align}
        \bigl(      A^\sigma\ketbra{b_2}{a_2}(A^\sigma)^\dagger
        \bigr)_{b_1,a_1}
        &=
        (A^\sigma)_{b_1b_2}
        \overline{(A^\sigma)_{a_1a_2}}.
        \notag
    \end{align}
    The right-factor formula is the same expansion with the two boundary bonds
    interchanged: the $(p,q)$ entry of $R_mR_m^\dagger$ is
    $\sum_\tau(A^\tau)_{p_2p_1}\overline{(A^\tau)_{q_2q_1}}$, which is the
    $(p_2,q_2)$ entry of $\E_A^{\,m}(\ketbra{p_1}{q_1})$.
\end{proof}

\begin{lemma}[RFP collapse of the operator-Schmidt Grams]
    \label{lem:rfp_schmidt_gram_collapse}
    \lean{MPSTensor.schmidtLeft_gram_eq_of_isTransferIdempotent,
        MPSTensor.schmidtRight_gram_eq_of_isTransferIdempotent}
    \leanok
    \uses{def:mps_transfer_idempotence, lem:schmidt_gram_transfer_power}
    Let $A$ be a renormalization fixed point. For every $L\ge1$ and
    $m\ge1$, the operator-Schmidt Grams of the $L$-block and of the
    $m$-site complement are the one-site Grams:
    \begin{align}
        L_L^\dagger L_L&=L_1^\dagger L_1,
        \notag\\
        R_mR_m^\dagger&=R_1R_1^\dagger.
        \notag
    \end{align}
\end{lemma}
\begin{proof}\leanok
    By definition, a renormalization fixed point satisfies
    $\E_A^{\,2}=\E_A$, so $\E_A^{\,k}=\E_A$ for
    every $k\ge1$. Applying Lemma~\ref{lem:schmidt_gram_transfer_power} to the
    left and right Gram entries gives the two displayed identities.
\end{proof}

\begin{lemma}[Pure block entropy as a bond-environment charpoly sum]
    \label{lem:pure_block_entropy_env_charpoly}
    \lean{MPSTensor.pureBlockEntropy_eq_env_charpoly}
    \leanok
    \uses{lem:schmidt_gram_transfer_power, lem:charpoly_root_entropy_cyclic_swap}
    Let $A$ be an MPS tensor, let $L\le N$, put $m=N-L$, and write
    $c=(\tr\rho^{(N)})^{-1}$. Then the pure block entropy is the
    charpoly-root entropy sum of the $D^2\times D^2$ bond environment:
    \begin{align}
        S_L^{(N)}(A)
        &=
        \sum_{\lambda\in
            \operatorname{roots}\chi_{c\,(R_mR_m^\dagger)(L_L^\dagger L_L)}}
            -\re(\lambda)\log\re(\lambda).
        \notag
    \end{align}
    The roots are counted with algebraic multiplicity.
\end{lemma}
\begin{proof}\leanok
    The reduced block state has the exact factorization
    \begin{align}
        \rho_L^{(N)}(A)
        &=
        c\,L_L(R_mR_m^\dagger)L_L^\dagger.
        \notag
    \end{align}
    Cyclic invariance of the charpoly-root entropy sum gives
    \begin{align}
        S(c\,L_L(R_mR_m^\dagger)L_L^\dagger)
        &=
        S(c\,(R_mR_m^\dagger)(L_L^\dagger L_L)),
        \notag
    \end{align}
    which is the displayed formula.
\end{proof}

\begin{theorem}[A renormalization fixed point saturates the area law]
    \label{thm:rfp-saturates-area-law}
    \lean{MPSTensor.isSAL_of_isTransferIdempotent}
    \leanok
    \uses{def:mps_transfer_idempotence, def:pure_sal,
        lem:pure_block_entropy_env_charpoly,
        lem:rfp_schmidt_gram_collapse}
    If an MPS tensor $A$ is a renormalization fixed point, then it saturates the
    area law: for every chain length $N$ and every block length
    $1\le L<\lfloor N/2\rfloor$, the pure-state block entropy is constant in the
    block size, $S_L^{(N)}(A)=S_{L+1}^{(N)}(A)$.
    The source statement is
    \cite[Proposition~3.14, lines~606--608]{Cirac2016MPDO_arXiv}.
    It assumes canonical form; the theorem above is the stronger pure-state
    implication obtained from transfer-map idempotence alone.
\end{theorem}
\begin{proof}\leanok
    \uses{lem:pure_block_entropy_env_charpoly, lem:rfp_schmidt_gram_collapse}
    Write $c=(\tr\rho^{(N)})^{-1}$. Lemma~\ref{lem:pure_block_entropy_env_charpoly}
    computes the entropy from the bond environment:
    \begin{align}
        S_L^{(N)}(A)
        &=
        \sum_{\lambda\in
            \operatorname{roots}\chi_{c\,(R_{N-L}R_{N-L}^\dagger)
                (L_L^\dagger L_L)}}
            -\re(\lambda)\log\re(\lambda).
        \notag
    \end{align}
    Since $1\le L<\lfloor N/2\rfloor$, the four lengths
    $L$, $L+1$, $N-L$, and $N-(L+1)$ are positive. The RFP Gram-collapse lemma
    gives
    \begin{align}
        L_L^\dagger L_L
        &=L_1^\dagger L_1,
        \notag\\
        L_{L+1}^\dagger L_{L+1}
        &=L_1^\dagger L_1,
        \notag\\
        R_{N-L}R_{N-L}^\dagger
        &=R_1R_1^\dagger,
        \notag\\
        R_{N-(L+1)}R_{N-(L+1)}^\dagger
        &=R_1R_1^\dagger.
        \notag
    \end{align}
    Therefore the two bond environments agree:
    \begin{align}
        (R_{N-L}R_{N-L}^\dagger)(L_L^\dagger L_L)
        &=
        (R_1R_1^\dagger)(L_1^\dagger L_1)
        \notag\\
        &=
        (R_{N-(L+1)}R_{N-(L+1)}^\dagger)
        (L_{L+1}^\dagger L_{L+1}).
        \notag
    \end{align}
    The scalar $c$ depends only on $N$, so the two charpoly-root entropy sums
    are equal. Hence $S_L^{(N)}(A)=S_{L+1}^{(N)}(A)$.
\end{proof}

\begin{theorem}[A renormalization fixed point has a constant block-entropy chain]
    \label{thm:rfp_constant_entropy}
    \lean{MPSTensor.pureBlockEntropy_eq_of_isTransferIdempotent}
    \leanok
    \uses{thm:rfp-saturates-area-law, lem:pure_sal_const}
    If an MPS tensor $A$ is a renormalization fixed point, then all block entropies
    in the range $1\le L\le\lfloor N/2\rfloor$ coincide:
    \begin{align}
        S_1^{(N)}(A)
        &=S_2^{(N)}(A)
         =\cdots
         =S_{\lfloor N/2\rfloor}^{(N)}(A).
        \notag
    \end{align}
    This is the explicit constant-entropy chain of the area-law saturation
    \cite[Definition~3.13, line~600]{Cirac2016MPDO_arXiv}, under the fixed-point
    hypothesis.
\end{theorem}
\begin{proof}\leanok
    \uses{thm:rfp-saturates-area-law, lem:pure_sal_const}
    The fixed point saturates the area law
    (Theorem~\ref{thm:rfp-saturates-area-law}), so by the telescoping of
    saturation (Lemma~\ref{lem:pure_sal_const}) the block entropies at any two
    lengths $1\le L,L'\le\lfloor N/2\rfloor$ agree:
    $S_L^{(N)}(A)=S_{L'}^{(N)}(A)$.
\end{proof}

\begin{lemma}[Operator-Schmidt rank bound]\label{lem:operator_schmidt_rank}
    \lean{MPSTensor.rank_reducedPureBlockState_le}
    \leanok
    For an MPS tensor $A$ of bond dimension $D$ and a block length $L\le N$, the
    reduced state $\rho_L^{(N)}(A)$ of the first $L$ spins has rank at most $D^2$.
\end{lemma}
\begin{proof}\leanok
    The reduced state is $\rho_L\propto WW^\dagger$, where the amplitude matrix
    \begin{align}
        W(u,w)
        &=
        \bra{u\,w}V^{(N)}(A)\rangle
        =
        \tr\!\bigl(      A^{u_1}\cdots A^{u_L}
          A^{w_1}\cdots A^{w_{N-L}}
        \bigr)
        \notag
    \end{align}
    factors through the two bond pairs cut by the block boundary. Writing the
    trace as a sum over the bond index pair
    $(a,b)\in\{1,\ldots,D\}^2$ exhibits $W=L'R'$ with $L'$ having $D^2$
    columns, so $\rank\rho_L\leq\rank W\leq D^2$.
\end{proof}

\begin{lemma}[The reduced pure block state is a density matrix]
    \label{lem:reduced_pure_block_state_density}
    \lean{MPSTensor.reducedPureBlockState_posSemidef, MPSTensor.reducedPureBlockState_trace}
    \leanok
    \uses{def:pure_sal}
    For an MPS tensor $A$ with $V^{(N)}(A)\neq0$, the reduced state
    $\rho_L^{(N)}(A)$ of the first $L$ spins is a density matrix: it is positive
    semidefinite and has unit trace.
\end{lemma}
\begin{proof}\leanok
    Positive semidefiniteness is preserved by the partial trace of the positive
    semidefinite normalized pure state, and the partial trace preserves the trace,
    which the normalization fixes to one.
\end{proof}

\begin{proposition}[Pure-state area-law upper bound]\label{prop:pure_area_upper}
    \lean{MPSTensor.pureBlockEntropy_le}
    \leanok
    \uses{lem:operator_schmidt_rank, lem:trace_one_positive_rank,
        thm:entropy_le_log_rank, lem:reduced_pure_block_state_density}
    Let $A$ be an MPS tensor of bond dimension $D\geq1$ with
    $V^{(N)}(A)\neq0$. For every block length $L\le N$,
    $S_L^{(N)}(A)\leq2\log D$, uniformly in $L$ and $N$. This is the area law of
    \cite[Section~3, line~599]{Cirac2016MPDO_arXiv}: the block entropy is bounded
    by a constant independent of the block size.
\end{proposition}
\begin{proof}\leanok
    \uses{lem:operator_schmidt_rank, lem:trace_one_positive_rank,
        thm:entropy_le_log_rank, lem:reduced_pure_block_state_density}
    The reduced state $\rho_L^{(N)}(A)$ is a density matrix
    (Lemma~\ref{lem:reduced_pure_block_state_density}), so its rank is positive
    by Lemma~\ref{lem:trace_one_positive_rank}. Its rank is at most $D^2$ by
    Lemma~\ref{lem:operator_schmidt_rank}. The rank bound on the entropy gives
    $S_L^{(N)}(A)\leq\log(D^2)=2\log D$.
\end{proof}

\begin{lemma}[Nonnegativity of the pure-state block entropy]
    \label{lem:pure_block_entropy_nonneg}
    \lean{MPSTensor.pureBlockEntropy_nonneg}
    \leanok
    \uses{def:pure_sal, thm:entropy_nonneg, lem:reduced_pure_block_state_density}
    For an MPS tensor $A$ with $V^{(N)}(A)\neq0$ and every block length
    $L\le N$, the pure-state block entropy is non-negative:
    $0\le S_L^{(N)}(A)$.
\end{lemma}
\begin{proof}\leanok
    \uses{lem:reduced_pure_block_state_density}
    The reduced state $\rho_L^{(N)}(A)$ is a density matrix
    (Lemma~\ref{lem:reduced_pure_block_state_density}), and the von Neumann
    entropy of a density matrix is non-negative.
\end{proof}

\begin{lemma}[Full-block pure-state entropy vanishes]
    \label{lem:pure_entropy_full}
    \lean{pureBlockEntropy_full}
    \leanok
    \uses{def:pure_sal}
    For an MPS tensor $A$ with $V^{(N)}(A)\neq0$, the entropy of the whole chain
    vanishes: $S_N^{(N)}(A)=0$.
\end{lemma}
\begin{proof}\leanok
    After reindexing the full block, the normalized operator is the rank-one
    pure state
    \begin{align}
        \frac{\ketbra{V^{(N)}(A)}{V^{(N)}(A)}}
        {\braket{V^{(N)}(A)}{V^{(N)}(A)}}.
        \notag
    \end{align}
    A pure state has von Neumann entropy zero.
\end{proof}

\begin{lemma}[Empty-block pure-state entropy vanishes]
    \label{lem:pure_block_entropy_empty}
    \lean{MPSTensor.pureBlockEntropy_empty}
    \leanok
    \uses{def:pure_sal, lem:pure_entropy_full, lem:schmidt_symmetry}
    For an MPS tensor $A$ with $V^{(N)}(A)\neq0$, the pure-state block entropy of
    the empty block vanishes: $S_0^{(N)}(A)=0$.
\end{lemma}
\begin{proof}\leanok
    \uses{lem:pure_entropy_full, lem:schmidt_symmetry}
    By the Schmidt symmetry $S_0^{(N)}(A)=S_N^{(N)}(A)$, and the entropy of the
    full system vanishes for a pure state
    (Lemma~\ref{lem:pure_entropy_full}).
\end{proof}

\begin{lemma}[Saturation telescopes to a constant block entropy]
    \label{lem:pure_sal_const}
    \lean{MPSTensor.pureBlockEntropy_eq_of_isSAL}
    \leanok
    \uses{def:pure_sal, lem:finite_interval_constancy}
    If an MPS tensor $A$ saturates the area law, then the block entropies
    coincide for all $1\le L,L'\le\lfloor N/2\rfloor$:
    $S_L^{(N)}(A)=S_{L'}^{(N)}(A)$.
\end{lemma}
\begin{proof}\leanok
    \uses{def:pure_sal, lem:finite_interval_constancy}
    The defining condition supplies
    $S_m^{(N)}(A)=S_{m+1}^{(N)}(A)$ throughout the interval
    $1\le m<\lfloor N/2\rfloor$. Apply
    Lemma~\ref{lem:finite_interval_constancy}.
\end{proof}

\begin{lemma}[Pure-state mutual information is twice the block entropy]
    \label{lem:mutual_info_doubled}
    \lean{mutualInfoChain_doubledTensor}
    \leanok
    \uses{lem:pure_entropy_full, lem:schmidt_symmetry, lem:blockEntropy_doubledTensor}
    For an MPS tensor $A$ with $V^{(N)}(A)\neq0$, the mutual information of the
    purification MPDO equals twice the pure-state block entropy:
    $I_L=2S_L^{(N)}(A)$.
\end{lemma}
\begin{proof}\leanok
    By definition, $I_L=S_L+S_{N-L}-S_N$ for the generated state. The global
    entropy of the pure state vanishes (Lemma~\ref{lem:pure_entropy_full}), and
    the Schmidt symmetry (Lemma~\ref{lem:schmidt_symmetry}) gives
    $S_{N-L}=S_L$, so $I_L=2S_L$.
\end{proof}

\begin{proposition}[Pure-state mutual-information bound]\label{prop:pure_mutual_info_bound}
    \lean{mutualInfoChain_doubledTensor_le}
    \leanok
    \uses{prop:pure_area_upper, lem:mutual_info_doubled}
    Let $A$ be an MPS tensor of bond dimension $D\geq1$ with
    $V^{(N)}(A)\neq0$. For the purification MPDO of $A$, whose generated
    operator is the pure state
    $\ketbra{V^{(N)}(A)}{V^{(N)}(A)}$, the mutual information between a block
    of $L$ spins and the rest satisfies $I_L\leq4\log D$, uniformly in $L$ and
    $N$.
\end{proposition}
\begin{proof}\leanok
    For a pure state, $I_L=2S_L^{(N)}(A)$, since the global entropy $S_N$
    vanishes and a block and its complement share the same Schmidt spectrum.
    The block-entropy bound $S_L^{(N)}(A)\leq2\log D$
    (Proposition~\ref{prop:pure_area_upper}) then gives $I_L\leq4\log D$.
\end{proof}
